\usepackage{csquotes} % Proporciona soporte para citas con formato adecuado, útil para trabajar con citas en diferentes idiomas

\usepackage[spanish,es-tabla]{babel} % Adapta el documento para el idioma español, incluyendo la adaptación de tablas y otros elementos lingüísticos

\usepackage[a4paper, margin=2cm, marginparwidth=3cm]{geometry} % Configura el tamaño del papel y los márgenes del documento

\usepackage{amsmath} % Ofrece soporte avanzado para la escritura de fórmulas matemáticas

\usepackage{graphicx} % Permite incluir y manipular imágenes en el documento

\usepackage{float} % Mejora la gestión de ubicaciones para figuras y tablas

\usepackage[colorlinks=true, allcolors=black]{hyperref} % Configura enlaces dentro del documento y en bibliografías con colores personalizables

\usepackage[export]{adjustbox} % Permite ajustar y recortar imágenes, y también usar herramientas como "frame" para enmarcar elementos

\usepackage[font=small, labelfont=bf]{caption} % Modifica la apariencia de los títulos de las figuras y tablas, cambiando el tamaño de la fuente y el formato de las etiquetas

\usepackage{subcaption} % Permite la creación de subfiguras y subtables dentro de figuras y tablas principales

\usepackage{physics} % Ofrece comandos para la notación matemática usada en física, como derivadas y operadores

\usepackage[table,xcdraw,x11names]{xcolor} % Permite el uso de colores en tablas y gráficos, facilitando la mejora de la apariencia de las tablas

\usepackage{booktabs} % Mejora el diseño de las tablas con líneas horizontales de diferentes grosores y estilos

\usepackage{todonotes} % Permite agregar notas y comentarios dentro del documento para su revisión

\usepackage{setspace} % Proporciona comandos para ajustar el espaciado entre líneas

\setuptodonotes{inline, size=\footnotesize, color=blue!20, bordercolor=none} % Configura cómo se mostrarán las notas de la biblioteca "todonotes" (en línea, tamaño pequeño, con color azul claro)

\usepackage{multicol} % Permite la creación de entornos multicolumna

\setlength\columnsep{18pt} % Especifica la separación entre columnas en entornos multicolumna

\usepackage{bigints} % Proporciona integrales grandes para su uso en ecuaciones matemáticas

\usepackage{tocloft} % Permite la personalización del índice (table of contents), listas de figuras y listas de tablas

\usepackage{wrapfig} % Permite envolver texto alrededor de figuras, ideal para mejorar la presentación de imágenes

\usepackage{fancyhdr} % Permite personalizar los encabezados y pies de página del documento

\usepackage{lastpage} % Permite hacer referencia a la última página del documento, útil para numeración de páginas

\pagestyle{fancy} % Activa el estilo de página personalizado definido por el paquete "fancyhdr"
\fancyhf{} % Limpia los encabezados y pies de página predeterminados

\cfoot{\thepage \hspace{1pt} de \pageref{LastPage}} % Define el formato del pie de página como "Página X de Y"

\usepackage{pdfpages} % Permite la inclusión de documentos PDF completos dentro del documento principal

\usepackage{bbold} % Proporciona una fuente para los números y caracteres en negrita doble (a menudo usada en matemáticas)

\usepackage{amssymb} % Proporciona símbolos matemáticos adicionales

\usepackage{enumitem} % Permite la personalización de listas enumeradas y con viñetas

\usepackage{amsmath,tikz} % Repite "amsmath" para fórmulas matemáticas avanzadas, y "tikz" para gráficos y diagramas vectoriales

\newcommand*\circled[1]{\tikz[baseline=(char.base)]{\node[shape=circle,draw,inner sep=1.5pt] (char) {#1};}} % Define un nuevo comando para crear números o caracteres dentro de círculos usando TikZ

\usepackage{circuitikz} % Permite dibujar circuitos eléctricos y electrónicos usando una sintaxis sencilla

\usepackage{tikz} % Permite crear gráficos y diagramas complejos en el documento usando una sintaxis versátil


\setlength{\headheight}{14pt} % Aumenta la altura del encabezado

\usepackage{fontawesome} % Permite insertar emojis y símbolos de Font Awesome


\usepackage{yfonts} %pone letras góticas con \textfrak

\usepackage[most]{tcolorbox}
\usepackage{lipsum} % Solo para texto de ejemplo

\newtcolorbox{infobox}[2][blue]{
    colback=#1!5!white,
    colframe=#1!75!black,
    title={#2},
    fonttitle=\bfseries,
    coltitle=black,
    boxrule=0.8pt,
    arc=4pt,
    left=6pt,
    right=6pt,
    top=6pt,
    bottom=6pt,
    breakable
}
\usepackage{cancel}
\newcommand{\colorcancel}[2]{\textcolor{#1}{\cancel{\textcolor{black}{#2}}}}

\usepackage{titlesec}

% Formato para párrafos como si fueran subtítulos normales
\titleformat{\paragraph}[block]{\normalfont\normalsize\bfseries}{\theparagraph}{1em}{}
\setcounter{secnumdepth}{4}
\setcounter{tocdepth}{4}

\usepackage{listings}

\lstdefinelanguage{custombash}{
    language=bash,
    morekeywords={conda, activate, mkdir, cd, cmake, make, git, sudo},
    sensitive=true
}

\lstdefinestyle{bashStyle}{
    language=custombash,
    basicstyle=\ttfamily\small,
    keywordstyle=\color{teal}\bfseries,
    commentstyle=\color{gray}\itshape,
    stringstyle=\color{orange},
    identifierstyle=\color{black},
    numbers=left,
    numberstyle=\tiny\color{gray},
    stepnumber=1,
    numbersep=10pt,
    backgroundcolor=\color{white},
    showspaces=false,
    showstringspaces=false,
    breaklines=true,
    frame=single,
    captionpos=b
}

\lstdefinestyle{pythonStyle}{
    language=Python,
    basicstyle=\ttfamily\small,
    keywordstyle=\color{blue}\bfseries,       % palabras clave
    commentstyle=\color{gray}\itshape,        % comentarios
    stringstyle=\color{orange},               % strings
    identifierstyle=\color{black},            % nombres de variables
    numbers=left,
    numberstyle=\tiny\color{gray},
    stepnumber=1,
    numbersep=10pt,
    backgroundcolor=\color{white},
    showspaces=false,
    showstringspaces=false,
    breaklines=true,
    frame=single,
    captionpos=b
}

\lstdefinestyle{cppStyle}{
    language=C++,
    basicstyle=\ttfamily\small,
    keywordstyle=\color{blue}\bfseries,
    commentstyle=\color{gray}\itshape,
    stringstyle=\color{red},
    numbers=left,
    numberstyle=\tiny\color{gray},
    stepnumber=1,
    numbersep=10pt,
    backgroundcolor=\color{white},
    showspaces=false,
    showstringspaces=false,
    breaklines=true,
    frame=single,
    captionpos=b
}

\usetikzlibrary{patterns}
\usepackage{bm}
\usetikzlibrary{patterns, decorations.markings, arrows.meta}
\usetikzlibrary{decorations.pathreplacing}
\usetikzlibrary{decorations.pathmorphing}
\definecolor{ForestGreen}{rgb}{0.13, 0.55, 0.13} % Define un color personalizado
\definecolor{Yellow}{HTML}{91760D} % Define un color personalizado
\definecolor{silver}{HTML}{DDDBD9} % Define un color personalizado